% =====================================================================
%  DESVIACION METODOLOGICA -- Diseno de analisis del experimento humano-LLM
%  Proyecto SIMPA -- Equipo AHMRV -- ISR-401 -- UTEQ -- 2026-2027 PPA
%  Compilacion: pdflatex osf_deviations.tex   (dos veces)
% =====================================================================
\documentclass[11pt,a4paper]{article}
\usepackage[utf8]{inputenc}
\usepackage[T1]{fontenc}
\usepackage{amsmath}
\usepackage[margin=2.4cm]{geometry}
\usepackage{longtable}
\usepackage{array}
\usepackage{xcolor}
\usepackage{colortbl}
\usepackage{enumitem}
\usepackage{fancyhdr}
\usepackage{lastpage}
\usepackage{titlesec}
\usepackage{url}
\usepackage[hidelinks]{hyperref}
\usepackage{caption}

\definecolor{uteqverde}{RGB}{0,102,51}
\definecolor{grisclaro}{RGB}{242,242,242}
\newcolumntype{L}[1]{>{\raggedright\arraybackslash}p{#1}}
\newcolumntype{C}[1]{>{\centering\arraybackslash}p{#1}}
\newcommand{\id}[1]{\texttt{#1}}
\renewcommand{\tablename}{Tabla}
\sloppy

\pagestyle{fancy}
\fancyhf{}
\lhead{\small Desviación metodológica --- SIMPA}
\rhead{\small Equipo AHMRV --- ISR-401}
\cfoot{\small Página \thepage\ de \pageref{LastPage}}
\renewcommand{\headrulewidth}{0.4pt}
\titleformat{\section}{\normalfont\large\bfseries\color{uteqverde}}{\thesection.}{0.6em}{}

\begin{document}

\begin{center}
{\large\bfseries UNIVERSIDAD TÉCNICA ESTATAL DE QUEVEDO\par}
{\normalsize Facultad de Ciencias de la Computación --- Carrera de Software\par}
\vspace{0.5cm}
\rule{\textwidth}{1pt}\par
\vspace{0.3cm}
{\LARGE\bfseries\color{uteqverde} DESVIACIÓN METODOLÓGICA\par}
\vspace{0.25cm}
{\large Diseño de análisis del experimento humano--LLM\\
\texttt{06\_Experimento/osf\_deviations.md}\par}
\vspace{0.3cm}
\rule{\textwidth}{1pt}\par
\end{center}

\vspace{0.3cm}
\begin{longtable}{|L{4.5cm}|L{10.3cm}|}
\hline
\rowcolor{grisclaro} \textbf{Proyecto} & SIMPA --- Equipo AHMRV --- ISR-401 --- UTEQ \\ \hline
\textbf{Responsable de esta declaración} & Denisses Huilcapi \\ \hline
\textbf{Verificador} & Allan Villafuerte \\ \hline
\textbf{Fecha de registro de este documento} & 11/09/2026 \\ \hline
\end{longtable}

\section{Cambio declarado}

\begin{longtable}{|L{3.0cm}|L{6.1cm}|L{6.1cm}|}
\hline
\rowcolor{uteqverde!15}
\textbf{} & \textbf{Protocolo original (OSF, \url{https://osf.io/4z35d/})} & \textbf{Diseño actualmente aplicable} \\ \hline
\endhead
Diseño &
Cuasi-experimento con medidas apareadas y evaluación a ciegas (\texttt{06\_Experimento/protocolo.tex}, línea 66) &
Grupos independientes, no apareados \\ \hline
Análisis principal &
Prueba $t$ apareada bilateral (o Wilcoxon con signo si no hay normalidad), corrección de Holm-Bonferroni &
Modelo mixto de vínculo acumulativo para respuesta ordinal, un modelo por cada una de las cinco dimensiones: \texttt{puntuación\_1-5 \textasciitilde{} origen + (1 | evaluador) + (1 | requisito)}, con \texttt{origen} como efecto fijo y \texttt{evaluador}/\texttt{requisito} como interceptos aleatorios cruzados. Enlace logístico de probabilidades proporcionales, salvo razón metodológica preespecificada para otro enlace. Implementación prevista: paquete \texttt{ordinal} de R (versión exacta a registrar en el momento de la ejecución) \\ \hline
Análisis de reserva &
--- &
Si el modelo mixto no converge: U de Mann-Whitney sobre la mediana por requisito entre evaluadoras, con $\delta$ de Cliff e IC 95\,\% como tamaño del efecto \\ \hline
Acuerdo entre evaluadoras (RQ2) &
$\kappa$ de Cohen (por pares) y $\kappa$ de Fleiss (conjunto) &
$\alpha$ de Krippendorff con métrica ordinal como medida principal; $\kappa$ ponderado cuadrático por pares como medida secundaria. Se descarta $\kappa$ nominal de Cohen/Fleiss porque, en una escala ordinal de 1 a 5, penaliza igual un desacuerdo de un punto que uno de cuatro \\ \hline
\end{longtable}

\section{Motivo del cambio}

El protocolo prerregistrado prevé contrastes apareados sobre veinticinco
pares. Al revisar el diseño antes de recoger puntuación alguna, se detectó
que no existe correspondencia uno a uno entre cada requisito humano y uno
generado por el LLM: ambos conjuntos proceden del mismo material fuente
(\texttt{ENTR-04}), pero no describen necesariamente las mismas funciones,
de modo que el emparejamiento carece de fundamento y el contraste apareado
sería inválido. Construir esa correspondencia a mano introduciría un juicio
subjetivo del propio equipo investigador, lo que comprometería la validez
de constructo del estudio.

Por esa razón se declara la siguiente desviación, sometida como adenda
prospectiva al registro \textbf{antes de recoger ninguna puntuación}:

\begin{enumerate}[label=\arabic*.]
  \item Los dos conjuntos se tratan como grupos independientes, no apareados.
  \item El análisis principal pasa a un modelo mixto ordinal (ver Sección 1),
        que respeta la naturaleza ordinal de la escala de cinco puntos y la
        estructura de evaluaciones repetidas (los mismos tres evaluadores
        puntúan ambos conjuntos) --- algo que un contraste sobre medias no
        captura correctamente.
  \item Como consecuencia directa del cambio de escala de análisis, el
        coeficiente de acuerdo entre evaluadoras también cambia: de
        $\kappa$ nominal a $\alpha$ de Krippendorff ordinal, por la misma
        razón de fondo (una métrica nominal no distingue la gravedad de un
        desacuerdo en una escala ordinal).
\end{enumerate}

Ninguna de estas tres decisiones se tomará ni se revisará en función de los
datos observados.

\section{Fecha en que se detectó}

El indicio más temprano verificable en el repositorio es el commit
\id{2a5412b} (``Incorporar manuscrito científico''), del \textbf{1 de
septiembre de 2026, 16:54:14 (-05:00)}, autoría de Denisses Huilcapi --- es
el primer commit que incorpora este razonamiento por escrito, en la sección
6.1 del manuscrito. Si el equipo identificó el problema antes en una
conversación no versionada, esa fecha anterior no cuenta con evidencia
documental verificable, por lo que esta declaración usa el 1 de septiembre
de 2026 como fecha de detección formal.

El protocolo original con diseño apareado se había cerrado y registrado
públicamente en OSF apenas el día anterior (commit \id{8bc6bde}, 31 de
agosto de 2026, ``Cerrar P8 con registro publico OSF del protocolo
experimental''). El experimento sigue sin ejecutarse a la fecha de este
documento (\texttt{06\_Experimento/readme.md}: \textit{``Protocolo diseñado
y registrado públicamente en OSF. Experimento no ejecutado.''}), de modo
que la desviación queda declarada con margen amplio respecto de cualquier
recolección real de puntuaciones.

\section{Impacto en el análisis planeado}

\begin{itemize}[leftmargin=1.1cm]
  \item \textbf{Prueba estadística original prevista:} prueba $t$ apareada
        bilateral (o Wilcoxon de rangos con signo si no hay normalidad),
        con corrección de Holm-Bonferroni sobre los cinco valores $p$
        (\texttt{protocolo.tex}, sección de plan de análisis).
  \item \textbf{Prueba estadística ahora aplicable:} modelo mixto ordinal
        (Sección 1), con la prueba U de Mann-Whitney + $\delta$ de Cliff
        como análisis de reserva si el modelo no converge.
  \item \textbf{Corrección por comparaciones múltiples:} se mantiene
        Holm-Bonferroni sobre los cinco contrastes del efecto de
        \texttt{origen}, uno por dimensión; se reportan valores $p$ crudos
        y corregidos, $\alpha = 0{,}05$. Esto no cambia respecto del
        protocolo original.
  \item \textbf{Acuerdo entre evaluadoras:} $\alpha$ de Krippendorff
        ordinal (principal) y $\kappa$ ponderado cuadrático por pares
        (secundario), en vez de $\kappa$ de Cohen/Fleiss nominal.
\end{itemize}

\subsection*{Potencia estadística --- recalculada para el nuevo diseño}

Con los mismos parámetros de referencia que fija el protocolo original
($\alpha = 0{,}05$ bilateral, potencia objetivo $1-\beta = 0{,}80$, tamaño
de efecto medio $d = 0{,}5$):

\begin{longtable}{|L{5.2cm}|C{4.7cm}|C{4.7cm}|}
\hline
\rowcolor{uteqverde!15}
\textbf{Diseño} & \textbf{Potencia con la muestra disponible} & \textbf{Tamaño requerido para potencia 0,80} \\ \hline
\endhead
Apareado (protocolo original) & --- & $\approx$ 34 pares (25 disponibles: insuficiente --- ya declarado en \texttt{07\_Datos/desviaciones.md}, D-01) \\ \hline
Grupos independientes (aproximación de sensibilidad, prueba $t$ de dos muestras) & $\approx$ 0,41 con 25 requisitos por grupo & $\approx$ 64 requisitos por grupo ($\approx$ 128 en total) \\ \hline
\end{longtable}

Este cálculo de sensibilidad ya está reportado en \texttt{manuscrito\_final.tex},
sección 6.3, y se reproduce aquí como referencia cruzada. \textbf{No
sustituye} una evaluación de potencia específica para el modelo mixto
ordinal, que exigiría simulación con las distribuciones y componentes de
varianza preespecificadas --- algo que el propio manuscrito reconoce como
pendiente.

\textbf{Conclusión de potencia:} el cambio de diseño agrava la limitación
de potencia ya declarada para el diseño apareado, no la mantiene igual. Se
reportará como amenaza a la validez de conclusión (Sección 7 del
manuscrito), priorizando el tamaño del efecto con su intervalo de confianza
al 95\,\% sobre la significación estadística. No se ampliará la muestra a
posteriori para corregir esta limitación.

\section{Relación con el protocolo original pre-registrado en OSF}

Este cambio se documenta como una \textbf{desviación}, no como una
modificación retroactiva del protocolo original. El registro OSF público
(\url{https://osf.io/4z35d/}) conserva intacta la versión del protocolo
presentada al momento del registro; esta declaración se somete como adenda
prospectiva, en la sección ``Deviations from pre-registration'' del
registro, antes de recoger ninguna puntuación.

\textbf{Responsable de actualizar la plataforma OSF:} Allan Villafuerte
(posee las credenciales de la cuenta del registro). Esta actualización en
osf.io es una acción pendiente independiente de subir este archivo al
repositorio --- ambas deben quedar hechas antes de que EXP-06 avance a la
sesión real de evaluación.

\section{Limitación reconocida}

El cambio de diseño de apareado a independiente reduce la potencia
estadística alcanzable respecto del diseño original, tal como muestra la
tabla de la Sección 4 (de $\approx$34 pares necesarios a $\approx$64 por
grupo, con una potencia observable de apenas $\approx$0,41 con la muestra
disponible). Esta limitación se reporta explícitamente en la sección de
Amenazas a la Validez del manuscrito (categoría: validez de conclusión),
siguiendo la misma práctica de transparencia ya aplicada en
\texttt{07\_Datos/desviaciones.md} (D-01 y D-02).

\vfill
\hrule
\vspace{0.2cm}
{\small Este documento en PDF se genera a partir de
\texttt{06\_Experimento/osf\_deviations.md}, fuente de referencia en caso de
discrepancia de formato.\par}

\end{document}
